\begin{abstract}
Open aviation trajectory data are now large enough to support learning-based forecasting at network scale, yet most existing aircraft trajectory studies still target short-horizon coordinate regression or terminal-area motion forecasting. These settings do not fully match the needs or the data quality of open flight archives, where long-horizon uncertainty, irregular coverage, and scalable evaluation are central concerns. We study multi-horizon regional aircraft trajectory forecasting on one month of open flight trajectories derived from EUROCONTROL's Aviation Data Repository for Research (ADRR)~\cite{eurocontrol_adrr_2025}. Given 30 minutes of flight history resampled at 5-minute intervals, the task is to predict future H3 region distributions at 10, 30, 60, and 120 minutes.

We formulate the problem as hierarchical coarse-to-fine H3 forecasting and propose a multiscale horizon-query Transformer that combines fine-scale temporal encoding, coarse pooled temporal encoding, and horizon-specific query attention over encoder states. The model outputs probabilistic regional forecasts together with candidate-region sets that can be evaluated by both accuracy and calibration. We evaluate on 20,712,602 forecast windows constructed from commercial flights operating in and over Europe, split by aircraft to reduce leakage across train, validation, and test sets.

Compared with a hierarchical GRU baseline, the proposed model improves macro top-1/top-3/top-5 accuracy from 0.6648/0.9076/0.9631 to 0.6827/0.9174/0.9685, reduces negative log-likelihood from 1.0214 to 0.9641, reduces Brier score from 0.4578 to 0.4338, and reduces expected calibration error from 0.0180 to 0.0093. Relative to a classical Markov regional transition baseline, gains are much larger, while candidate regions also become substantially smaller. An H3-neighbor refinement postprocessor improves candidate recall but degrades calibration, making it better suited as an ablation than as the default inference path. These results support regional probabilistic forecasting as a scalable alternative to long-horizon point prediction on open aviation trajectory data.
\end{abstract}
